\documentclass[12pt,a4paper]{config/myConfig}
\usepackage{silence}
\WarningFilter{latex}{You have requested}
\usepackage{config/myDependencies}
\usepackage{style/myEnv}
\usepackage{style/myCommands}
\usepackage{style/myColors}

\addbibresource{bibliography.bib}  % your .bib file

%Number of the practica for this assigment:
\newcommand{\AssigmentNumber}{}%<- Nombre de la Entrega
\newcommand{\AssigmentDate}{} %<- Fecha de Entrega
%------------------
%If any computer language used:
\newcommand{\Rcode}{R - versión 3.4.4}
\newcommand{\Pycode}{Python - versión 3.13}
%----------------------
%Palabras Clave para el abstract del informe:
\newcommand{\PalabrasClave}{\small \it
..,$\ $.. ,$\ $ ..}

\title{\Large $\ $\\ \bf } %<-Titulo del informe

\author
\info

\abstract{\PalabrasClave
\normalsize
\\[17pt]
{\bf Abstract.} 
Some Text...
}

%Table specifications:
\sisetup{
  separate-uncertainty = true,
  table-number-alignment = center,
  detect-weight = true,
  detect-family = true
}

\begin{document}
\thispagestyle{myheadings}
\pagestyle{myheadings}
\markright{\tt \AssigmentNumber|    \AssigmentDate}%check the date



\section{Marco Conceptual}
\label{sec:MARCO-CONCEPTUAL}

\subsection{Antecedentes}
\label{sec:ANTECEDENTES}

\section{Introducción}
\label{sec:INTRODUCCION}

\subsection{Planteamiento del Problema}
\label{sec:PLANTAMIENTO-PROBLEMA}

\subsection{Supuestos de la Investigación}
\label{sec:SUPUESTOS}

\begin{supuestos}
  \item \label{su:uno} ...
  \item \label{su:dos} ...
\end{supuestos}

\section{Hipótesis y Objetivos}
\label{sec:HIPOTESIS-OBJECTIVOS}

\subsection{Hipótesis}
\label{sec:HIPOTESIS}


\begin{hipotesis}
  {} %<- La hiptesis escrita en palabras.
  {} %<- hipotesis nula
  {} %<- hipotesis alternativa
  {} %<- Modelo plantiado para probar la hipotesis
  {} %<- Suposuciones para el Modelo como tal.
\end{hipotesis}


\subsection{Objetivos}
\label{sec:OBJECTIVOS}

\begin{objectives}
  \item \label{obj:A} ...
  \item \label{obj:B} ... 
\end{objectives}

\section{Metodología}
\label{sec:METODOLOGIA}

\subsection{Presentación de Montaje}
\label{sec:MONTAJE}

\begin{FigureLab}{Experimento Montado 1}{}{fig:Montaje 1}
      \includegraphics[width=0.5\linewidth]{}
\end{FigureLab}


%En el texto: Como se muestra en la \cref{fig:Montaje 1}, \FigRef{fig:Montaje 1} 


\subsection{Materiales}
\label{sec:MATERIALES}

\begin{TableLab}{Tabla de Materiales}{}{tab:Materiales}
    \rowcolors{2}{gray!10}{white}
    \newcolumntype{Y}{>{\raggedright\arraybackslash}X}
    \begin{tabularx}{\textwidth}{@{} l l >{\hsize=0.6\hsize}Y  >{\raggedright\arraybackslash}X @{}}
        \toprule
        \rowcolor{lightOrange}
        \textbf{Material} & \textbf{Cantidad} & \textbf{Especificación} \\
        \midrule
         &  &  \\
         &  &  \\
        \bottomrule
    \end{tabularx}
\end{TableLab}

%\cref{tab:Materiales} or \TabRef{tab:Materiales}


\subsection{Metodos}
\label{sec:METODOS}


\section{Ecuaciones}
\label{sec:ECUACIONES}


%Example of an equation
\begin{EquationLab}
{} % <- ecuacion con simbologia matematica
{
\begin{itemize}[label=$\circ$, left=0pt]
    \item $ $: ...
\end{itemize}
} % <- Indicacion de cada simbolo puesto en la ecuacion.
\end{EquationLab}

%\cref{eq:1}

\section{Datos}
\label{sec:DATOS}

%The physical parameters or Constants that where used for the experiment:
\begin{itemize}[label=$\rightarrow$, left=0pt]

  
  %\item para1 : num $\left( unit \right)$\hfill(\cite{})
  %Parameter, Value, Unit, Reference where obtained.
\end{itemize}


\section{Análisis}
\label{sec:RESULTADOS}


\subsection{Cálculos}
\label{sec:CALCULOS}

%Here we test numerically each of the ecuations proposed to indicate how where each ecuations measured.
\begin{calculoLab}
    { } % <- Name of the ecuation 
    {\cref{eq:}} % <- indication via hyperlink on which equation used.
    {
     \begin{array}{rcl}
        &&   
    \end{array}
    } % <- Arthimetically demostration.
\end{calculoLab}


\subsection{Resultados}
\label{sec:RESULTADOS}

%Example of a table:
%para: Title of Table,  if Code is used i.e Rcode , if different hyperreference is used (else, use \cref{tab:}.
\begin{TableLab}{}{}{}
  \begin{tabular}
    
  \end{tabular}
\end{TableLab}






\section{Discusión de Resultados}
\label{sec:DISCUSSION}
    





%Como se ve en la \cref{tab:II},
%\cref{hyp:A}
%\hyperref[obj:A]{Objetivo A}

%for apa7 reference use: \parencite{smith2019study}  (author, date)
%to use in text citation use: \textcite{garcia2020efecto}. author (date)

%To reference an objective do~\hyperref[obj:A]{Objetivo A}
%To reference a equation do~\autoref{eq:equationLabel} 

%To reference a table use ~\autoref{tab:TableNO} this prints the label and a hyperlink to that table


%To reference a figure use ~\autoref{fig:figNo} this prints the figure number and the hyperlink
\section{Conclusiones}
\label{sec:CONCLUSIONES}

\begin{itemize}
    \item ..
    \item ..
    \item ..
\end{itemize}




\footnotesize
\printbibliography
\normalsize

\section*{Anexo}
\label{sec:Anexo}

%param: Title, Code, Hyperlink
\begin{FigureAnexo}{}{}{}
\includegraphics[width=0.8\linewidth]{}
\end{FigureAnexo}

\end{document}
 